\section{Compact examples: forbidden, mixed and obstructed deformations}
\label{sec:examples}

Three compact configurations exhibit distinct consequences of the charge
map: a local Higgs flow that disappears, a mixed transition that exists only
through cooperation, and a large system with many partial directions but no
simultaneous all-sector trajectory. These effects depend on the actual kernel,
not only its dimension.

\subsection[A single E2 direction that disappears globally]{A single $E_2$ direction that disappears globally}
\label{subsec:x7}

The smallest example, denoted $X_7$, is a compact anticanonical hypersurface
with one isolated cone over $S_7$ and no other singularities. Locally, the
$E_2$ theory has the reduced Higgs coordinate $\beta$ of
Section~\ref{subsec:e2-sector}. In the convention
\eqref{eq:polar-convention}, it is a $\Delta_7$-regular hypersurface in
$\mathbb P_{\Delta_7}$, where
\begin{equation}
\begin{split}
 \Delta_7=\operatorname{conv}\{&v_0=(1,0,0,0),v_1=(0,1,0,0),
 v_2=(0,0,1,0),v_3=(-6,-4,-1,0),\\
 &v_4=(0,0,0,1),v_5=(-6,-4,0,-1),
 v_6=(-3,-2,1,-1)\}.
\end{split}
 \label{eq:x7-polytope}
\end{equation}
This reflexive polytope has fifteen unimodular triangular two-faces and one
unit-edge pentagon
$G_7=\operatorname{conv}\{v_2,v_3,v_4,v_5,v_6\}$ with the unique interior
lattice point $p_7=(-3,-2,0,0)$; its dual edge has length one. The three
remaining boundary lattice points of $\Delta_7$ lie in facet interiors, so
their divisors miss a generic hypersurface. We fix a
maximal projective crepant partial (MPCP) subdivision whose restriction to
$G_7$ is the five-triangle star at
$p_7$. All other two-faces are already unimodular, so this is the only local
resolution choice entering the root calculation. Writing $D_i$ for the ray
divisor of $v_i$, the reduced $A_1$ root pairs with the ambient divisor
classes as
\begin{equation}
 b_{7,\beta}=-D_2-D_3+D_4+D_5.
 \label{eq:x7-root-row}
\end{equation}
The row is nonzero and the reduced local branch is one-dimensional. Hence
the root class has nonzero image in $H_2(\widehat X_7;\QQ)$, and therefore
\begin{equation}
 K_{\mathrm{red}}(X_7)=0.
 \label{eq:x7-kernel}
\end{equation}
The corresponding compact vector gauges the $SU(2)$ Cartan with nonzero
coefficient. Its complex moment-map equation sets $\beta=0$.

Equation \eqref{eq:x7-kernel} has two consequences at different logical
levels. Theorem~\ref{thm:first-order-image} says that no global first-order
complex deformation moves the reduced local $E_2$ coordinate. The
all-orders Theorem~\ref{thm:all-orders-necessity} then excludes a
one-parameter smoothing whose first nonzero local term appears at higher
order. Thus a genuine local Higgs branch of the isolated fixed point is absent
from the compact deformation space.

This example is intentionally elementary. A single local sector supplies
nothing with which its nontrivial root class can cancel. It provides the clean
negative control for the cooperative example below.

\subsection[A mixed smoothing requiring simultaneous local deformations]{A mixed smoothing requiring simultaneous local deformations}
\label{subsec:x9}

The main positive example $X_9$ contains two nodes and one cone over
$S_7$. Order the reduced invariant coordinates as
\begin{equation}
 x=(u_1,u_2,\beta)^{\mathsf T},
 \label{eq:x9-coordinates}
\end{equation}
where $u_1,u_2$ are conifold moment maps and $\beta$ is the $E_2$ Cartan
moment map. No nonzero relation among the two nodal curve classes alone has a
nonzero coefficient on either curve.
The $E_2$ root likewise cannot move by itself, as the preceding example
illustrates.

The underlying reflexive polytope is
\begin{equation}
\begin{split}
 \Delta_9=\operatorname{conv}\{&v_0=(1,0,0,0),v_1=(0,1,0,0),
 v_2=(-1,-1,0,0),\\
 &v_3=(0,0,1,0),v_4=(0,0,0,1),v_5=(-4,-2,-1,0),\\
 &v_6=(-4,-2,0,-1),v_7=(3,1,1,1),v_8=(2,1,1,1)\}.
\end{split}
 \label{eq:x9-polytope}
\end{equation}
Its singular two-faces are the squares
$\{v_2,v_3,v_6,v_7\}$ and $\{v_2,v_4,v_5,v_7\}$ and the pentagon
$\{v_3,v_4,v_5,v_6,v_8\}$. Star subdivision at the pentagon's interior
point $p=(-2,-1,0,0)$ gives $E=D_p\cong S_7$. The only other boundary
lattice point, $(-1,0,0,0)$, lies in a facet interior, and its divisor misses
a generic anticanonical hypersurface. We fix an MPCP completion using the
square diagonals $(v_6,v_7)$ and $(v_2,v_4)$, in the order just listed. With
the sign convention that the chosen diagonal has coefficient $-1$, the two
nodal circuit rows on the ten rays $(v_0,\ldots,v_8,p)$ are
\begin{equation}
 v_2+v_3-v_6-v_7,
 \qquad
 -v_2-v_4+v_5+v_7.
 \label{eq:x9-node-circuits}
\end{equation}
The remaining choices in the MPCP completion lie over smooth strata and do
not alter these circuit rows or the star of $p$. Principal-divisor relations eliminate
$D_0,D_1,D_3,D_4$, giving the divisor basis below.

After removing a toric divisor that does not meet the generic hypersurface, a
saturated basis of the toric divisor quotient is
\begin{equation}
 \mathcal D=(D_2,D_5,D_6,D_7,D_8,E),
 \label{eq:x9-divisor-basis}
\end{equation}
where $E\cong S_7$ is the exceptional divisor. This is a complete basis of
$H^2(\widehat X_9;\QQ)$. In this basis, the compact gauging matrix is
\begin{equation}
 Q_{\mathcal D}
 =\begin{pmatrix}
  1&-1& 0\\
  0& 1& 1\\
 -1& 0&-1\\
 -1& 1& 0\\
  0& 0& 0\\
  0& 0& 0
 \end{pmatrix}.
 \label{eq:x9-full-matrix}
\end{equation}
An integral unimodular change of the six compact vector directions isolates
two effective vectors,
\begin{equation}
 V_1=-D_6,
 \qquad
 V_2=D_5,
 \label{eq:x9-effective-vectors}
\end{equation}
and brings \eqref{eq:x9-full-matrix} to
\begin{equation}
 Q_{X_9}^{\mathrm{eff}}
 =\begin{pmatrix}1&0&1\\0&1&1\end{pmatrix}.
 \label{eq:x9-effective-matrix}
\end{equation}
Both nonzero Smith invariants are one in the printed toric lattice. The two
complex moment-map equations are
\begin{equation}
 u_1+\beta=0,
 \qquad
 u_2+\beta=0,
 \label{eq:x9-moment-maps}
\end{equation}
so
\begin{equation}
 \ker Q_{X_9}^{\mathrm{eff}}
 =\CC(-1,-1,1).
\label{eq:x9-kernel}
\end{equation}

The cooperative structure is displayed in
Figure~\ref{fig:x9-cooperation}.
\begin{figure}[t]
\centering
\begin{adjustbox}{max width=\textwidth,center}
\begin{tikzpicture}[
  node distance=7mm and 7mm,
  sector/.style={draw,rounded corners,fill=gray!7,align=center,
    minimum height=12mm,text width=0.19\textwidth,font=\small},
  equation/.style={draw,very thick,rounded corners,align=center,
    minimum height=13mm,text width=0.31\textwidth,font=\small},
  endpoint/.style={draw,double,rounded corners,align=center,
    minimum height=13mm,text width=0.27\textwidth,font=\small},
  note/.style={draw,dashed,rounded corners,align=left,
    text width=0.27\textwidth,font=\scriptsize},
  arrow/.style={-{Latex[length=2.2mm]},thick}
]
\node[sector] (h1) {conifold hyper 1\\$u_1=h_1\widetilde h_1$};
\node[sector,right=of h1] (h2) {conifold hyper 2\\$u_2=h_2\widetilde h_2$};
\node[sector,right=of h2] (e2) {$E_2$ fixed point\\$\beta=\mu^{\CC}_{SU(2)}$};

\node[equation,below=11mm of h2] (mm)
  {$Q^{\mathrm{eff}}_{X_9}=\begin{psmallmatrix}1&0&1\\0&1&1\end{psmallmatrix}$\\[1mm]
   $u_1+\beta=0$\,,\quad $u_2+\beta=0$};
\node[endpoint,below=9mm of mm] (smooth)
  {primitive mixed trajectory\\
   $(u_1,u_2,\beta)=t(-1,-1,1)$};
\node[note,right=11mm of mm] (vectors)
  {$V_1,V_2$: flavor vectors Higgsed by the rank-two gauging.\\[1mm]
   $E$: neutral $E_2$ Coulomb vector lost when the exceptional divisor is
   smoothed; it is not a third gauging.};

\draw[arrow] (h1.south) -- ++(0,-5mm) -| ([xshift=-6mm]mm.north);
\draw[arrow] (h2.south) -- (mm.north);
\draw[arrow] (e2.south) -- ++(0,-5mm) -| ([xshift=6mm]mm.north);
\draw[arrow] (mm.south) -- node[right,font=\scriptsize] {$\ker Q_{X_9}^{\mathrm{eff}}$} (smooth.north);
\draw[arrow,dashed] (mm.east) -- (vectors.west);
\end{tikzpicture}
\end{adjustbox}
\caption{Cooperative Higgsing in $X_9$. No nonzero trajectory is supported on
the two conifold sectors or on the $E_2$ sector alone. The unique global
direction activates all three, while the protected vector-multiplet count
separates the two Higgsed flavor vectors from the neutral SCFT Coulomb vector.}
\label{fig:x9-cooperation}
\end{figure}


No nonzero vector in \eqref{eq:x9-kernel} is supported only on the nodes or
only on the $E_2$ sector. The two free-hypermultiplet sectors and the
interacting sector can move only together. Conversely, an explicit
one-parameter Calabi--Yau family realizes the class
\begin{equation}
 (u_1,u_2,\beta)=t(-1,-1,1)
 \label{eq:x9-family-direction}
\end{equation}
and has smooth general fiber by Theorem~11.8 of
\cite{Wuebben:2026smoothing}. More explicitly, take the facet
$F=\{v_2,v_3,v_4,v_5,v_6,v_7,v_8\}$ and deformation degree
$R=(0,1,-1,-1)$. On the ordered edges
\begin{equation}
 (v_2v_5,v_2v_6,v_2v_7,v_3v_6,v_3v_7,v_3v_8,
 v_4v_5,v_4v_7,v_4v_8,v_5v_6,v_7v_8)
 \label{eq:x9-facet-edges}
\end{equation}
the edge-dilation vector is
$\tau=(0,0,1,1,0,0,1,0,0,0,0)$. This is the compatible Minkowski
decomposition used in the Ilten--Vollmert construction
\cite{Altmann:1995,Ilten:2012}: it restricts to a
segment-plus-segment decomposition on each square and to a
segment-plus-unimodular-triangle decomposition on the pentagon. These data
define the cited toric deformation and its relative anticanonical family. The
kernel is
therefore not only a first-order flat direction: it integrates to a mixed
extremal transition, and it is cooperative in the sense defined in the
Introduction.

There is a useful spectrum check. Exact lattice-point and
Mayer--Vietoris calculations give
\begin{equation}
 \begin{array}{c|ccc|cc}
 &h^{1,1}&h^{2,1}&\chi&n_V&n_H^{\mathrm{neutral}}\\
 \hline
 \widehat X_9&6&122&-232&5&123\\
 X_{9,t}&3&123&-240&2&124.
 \end{array}
 \label{eq:x9-spectrum-table}
\end{equation}
For a smooth compact Calabi--Yau threefold in M-theory,
$n_V=h^{1,1}-1$ and $n_H^{\mathrm{neutral}}=h^{2,1}+1$. The two effective
gaugings in \eqref{eq:x9-effective-matrix} Higgs two vector directions and
remove two quaternionic directions from the product of local Higgs sectors.
One quaternionic modulus remains, in agreement with the increase of
$n_H^{\mathrm{neutral}}$ in \eqref{eq:x9-spectrum-table}.

The transition loses three, not two, massless vector multiplets. The third has
a different origin. The exceptional-divisor vector $E$ is neutral under the
$E_2$ root moment map because
\begin{equation}
 E|_E=K_E,
 \qquad
 \beta\in K_E^{\perp}.
 \label{eq:x9-neutral-exceptional}
\end{equation}
It is the rank-one $E_2$ Coulomb direction and disappears when the entire
surface is replaced by the smoothing. It is not a third flavor gauging. This
separation between Higgsed vectors and the lost SCFT Coulomb vector is needed
for the protected spectrum to agree with the topology.

The integrality statement above is made in the saturated toric divisor
lattice. A possible finite index between that lattice and the full integral
cohomology lattice would affect global forms or discrete remnants, but not the
rank, the complex kernel or the multiplet count. We also do not claim the full
massive spectrum at the interacting $E_2$ point.

\subsection{Thirty locally smoothable sectors with no common trajectory}
\label{subsec:xcirc}

The large example $X^{\circ}$ contains
\begin{equation}
 26\text{ nodes},
 \qquad 2\text{ }E_2\text{ points},
 \qquad 2\text{ }E_3\text{ points}.
 \label{eq:xcirc-content}
\end{equation}
Every singular germ admits a local smoothing. At the $i$th $E_2$ point, let
$\alpha_i$ denote the nilpotent scheme-theoretic tangent coordinate and
$\beta_i$ the reduced smoothing coordinate. At the
$j$th $E_3$ point, let $L_j$ be the $A_1$ coordinate and
$(H_{j1},H_{j2})$ the $A_2$ coordinates. Record a profile by the ordered pair
of component labels selected at the two $E_3$ points. The four reduced
profiles are
\begin{equation}
 (2,2),\quad (2,3),\quad (3,2),\quad (3,3),
 \label{eq:xcirc-profiles}
\end{equation}
where $2$ and $3$ are subscripts naming the selected simple factor:
$2$ selects the $\mathfrak{sl}_2$ line and $3$ the $\mathfrak{sl}_3$ plane.

Appendix~\ref{app:xcirc-input} gives the defining fan polytope $\Pi$, all singular
two-faces, a strict convex height function fixing the crepant resolution, and
the complete sparse charge matrix. Thus $X^{\circ}$ denotes a specified
$\Pi$-regular anticanonical hypersurface rather than only an abstract
singularity inventory.

In an integral divisor frame the full local-to-global matrix has size
$36\times26$ and rank $21$; the physical charge matrix is its transpose.
Restricting to the four profiles gives the exact data in
Table~\ref{tab:xcirc-profiles}.

\begin{table}[t]
\centering
\small
\begin{adjustbox}{max width=\textwidth,center}
\begin{tabular}{@{}ccccc@{}}
\toprule
profile & $\dim_{\CC}\cB_\sigma$ & $\rank Q_\sigma$ & $\dim_{\CC}\ker Q_\sigma$ & forced reduced coordinates \\
\midrule
$(2,2)$ & $30$ & $21$ & $9$ & $u_9,u_{11},\beta_1,\beta_2$ \\
$(2,3)$ & $31$ & $21$ & $10$ & $u_9,u_{11},\beta_1,\beta_2,L_1$ \\
$(3,2)$ & $31$ & $21$ & $10$ & $u_9,u_{11},\beta_1,\beta_2,L_2$ \\
$(3,3)$ & $32$ & $20$ & $12$ & $u_9,u_{11},\beta_1,\beta_2$ \\
\bottomrule
\end{tabular}
\end{adjustbox}
\caption{The four branch-restricted compact charge maps for $X^{\circ}$.
In a profile $(\sigma_1,\sigma_2)$, $\sigma_j=2$ selects the
$\mathfrak{sl}_2$ line and $\sigma_j=3$ the $\mathfrak{sl}_3$ plane at the
$j$th $E_3$ point; $\alpha_i$ is the nilpotent scheme-theoretic tangent
coordinate and $\beta_i$ the reduced smoothing coordinate at the $i$th $E_2$
point, while $L_j$ is the $\mathfrak{sl}_2$ line coordinate
at the $j$th $E_3$ point. Every nonzero Smith invariant is one in the printed
divisor lattice.}
\label{tab:xcirc-profiles}
\end{table}

The forced-zero conclusions are witnessed by individual compact vectors, not
only by a rank calculation. Five primitive integral vector combinations have
complex moment maps
\begin{align}
 \mu_{V_9}^{\CC}&=u_9-\alpha_1+\alpha_2,&
 \mu_{V_{11}}^{\CC}&=u_{11}-\alpha_1+\alpha_2,
 \nonumber\\
 \mu_{V_{\beta_1}}^{\CC}&=\beta_1,&
 \mu_{V_{\beta_2}}^{\CC}&=\alpha_2+\beta_2,
 \nonumber\\
 \mu_{V_{A_1}}^{\CC}&=L_1+L_2.
 \label{eq:xcirc-isolating-moments}
\end{align}
At an ordinary complex point of the $E_2$ Higgs scheme,
$\alpha_1=\alpha_2=0$. The first four equations in
\eqref{eq:xcirc-isolating-moments} therefore force
\begin{equation}
 u_9=u_{11}=\beta_1=\beta_2=0
 \label{eq:xcirc-universal-zeros}
\end{equation}
on every reduced profile. The fifth equation correlates $L_1$ and $L_2$ on
$(2,2)$, isolates the remaining line coordinate on $(2,3)$ or $(3,2)$, and
becomes inactive on $(3,3)$. This explains the rank drop in the last row of
Table~\ref{tab:xcirc-profiles}.

The correct conclusion is not that $X^{\circ}$ has no Higgs directions. Its
projected kernels have dimensions between nine and twelve. Rather, no profile
contains a point at which all thirty local sectors have nonzero smoothing
coordinates. The all-orders result is stronger than this first-order
observation. Any one-parameter smoothing of $X^{\circ}$ would smooth the
nodes $N_9$ and $N_{11}$ and both $E_2$ points, whatever components it
selects at the two $E_3$ points. Theorem~\ref{thm:all-orders-necessity} and
Remark~\ref{rem:full-support} then require a class in the kernel of the
induced profile with nonzero projection to the coordinates in
\eqref{eq:xcirc-universal-zeros}, and Table~\ref{tab:xcirc-profiles}
excludes such a class on every profile. Consequently no analytic
one-parameter family smooths all thirty germs, even if its local parameters
begin at different orders. This is Theorem~6.5 of
\cite{Wuebben:2026obstruction}, whose Corollary~6.7 also excludes every
formal arc meeting all thirty local smoothing loci.

\subsection{Nonzero deformations constrained to a discriminant}
\label{subsec:discriminant-obstruction}

There is a second obstruction mechanism: an allowed $E_3$ deformation
can be nonzero but confined to a discriminant line. We use the compact
example of \cite[Theorems~5.1 and~5.3]{Wuebben:2026smoothing}, with thirty
nodes and two degree-six cone points. This is a different threefold from
$X^{\circ}$, which has twenty-six nodes and four cone points. The cited
construction specifies its reflexive fan polytope, projective crepant
resolution and full divisor--curve matrix. Here we record the three
divisor identities that explain its obstruction as compact gauging.

Let $n_0,\ldots,n_{29}$ be its nodal coordinates. At cone $i=0,1$,
write $\lambda_i$ for the $A_1$ coordinate and $(a_i,b_i)$ for the
coefficients of $A_i=e_3-e_1$ and $B_i=e_2-e_3$ in its $A_2$ plane.
The three plane coefficients are
\begin{equation}
 a_iA_i+b_iB_i=-a_i e_1+b_i e_2+(a_i-b_i)e_3,
 \qquad a_ib_i(a_i-b_i)\ne0
 \label{eq:discriminant-example-plane}
\end{equation}
for a smoothing direction. Three integral compact divisor combinations
give the moment-map equations
\begin{equation}
 \begin{aligned}
 n_{13}-a_0+b_0-a_1+b_1&=0,\\
 b_0+b_1&=0,\\
 n_{16}+\lambda_0+\lambda_1&=0.
 \end{aligned}
 \label{eq:discriminant-example-gaugings}
\end{equation}
They hold on the full charge kernel before a profile is chosen.
Denote the line and plane by $L$ and $P$. On $LL$ the first equation
forces $n_{13}=0$, while on $PP$ the third forces $n_{16}=0$.
On $LP$ the second forces $b_1=0$; on $PL$ it forces $b_0=0$.
Every profile therefore fails the exact smoothing criterion.

The mixed profiles nevertheless contain relations nonzero at every
singular point. For example the cited $LP$ relation has
$(\lambda_0,a_1,b_1)=(3,1,0)$ and all thirty nodal coefficients nonzero.
Its $E_3$ plane deformation is nonzero, but its middle coefficient in
\eqref{eq:discriminant-example-plane} vanishes. Thus a nonzero invariant
deformation at each sector does not imply simultaneous smoothing.
The selected bases are smooth of dimensions $34,34,34,35$ on
$LL,LP,PL,PP$, respectively, and all their tangents integrate
\cite[Section~5.3]{Wuebben:2026smoothing}. The failure is to leave the
local discriminants, not a failure of selected-base integrability. This
does not exclude partial Higgs directions or identify the entire physical
vacuum space with a deformation base.

\subsection{Finite classification for polytopes with at most nine vertices}
\label{subsec:census}

The three examples are part of a finite, reproducible classification. Among
the $77$ admissible reflexive four-polytopes with at most nine vertices in
the Kreuzer--Skarke list \cite{Kreuzer:2000} and a degree-seven cone point,
the all-orders necessity condition obstructs $76$;
$67$ of these have a single singular point, while nine have one node in
addition to the $E_2$ point. The remaining case is $X_9$, whose explicit mixed
family is smooth. Thus cooperative Higgsing is
the unique smoothable outcome in this precisely stated bounded class
by Proposition~9.2 and Corollary~6.7 of \cite{Wuebben:2026obstruction} and
Theorems~11.8 and~10.3 of \cite{Wuebben:2026smoothing}.

This census shows that the compact constraint is systematic rather than an
accident of one matrix. It is not a statistical sample of all compact
Calabi--Yau threefolds, and no universality claim is inferred from the ratio
$76/77$.
